% Physical source leaf 051, printed46; article55 proofII continuation.
\SourceLeaf{51}
Metietur itaque $t$ ipsum $\frac{p-1}{a}$; quare productum $ABC$ etc. etiam ad potestatem ${\frac{p-1}{a}}^{\mathrm{tam}}$ elevatum unitati erit congruum (art. 46). Sed perspicuum est singulos $B,C$, etc. (exemto ipso $A$) ad potestatem ${\frac{p-1}{a}}^{\mathrm{tam}}$ elevatos unitati congruos fieri, quum exponentes $b^\beta,c^\gamma$ etc. ad quos singuli pertinent ipsum $\frac{p-1}{a}$ metiantur. Hinc erit
\[A^{\frac{p-1}{a}}B^{\frac{p-1}{a}}C^{\frac{p-1}{a}}\text{ etc.}\equiv A^{\frac{p-1}{a}}\equiv1.\]
Unde sequitur exponentem ad quem $A$ pertinet ipsum $\frac{p-1}{a}$ metiri debere (art. 48), \emph{i. e.} $\frac{p-1}{a^{\alpha+1}}$ esse integrum; at $\frac{p-1}{a^{\alpha+1}}=\frac{b^\beta c^\gamma\text{ etc.}}{a}$ integer esse nequit (art 15). Unde tandem concludere oportet, suppositionem nostram consistere non posse, \emph{i. e.} productum $ABC$ etc. revera ad exponentem $p-1$ pertinere. \emph{Q. E. D.}

Demonstratio posterior priori aliquantulum prolixior esse videtur, prior contra posteriori minus directa.

\Article{56}
Hoc theorema insigne exemplum suppeditat, quanta circumspectione in theoria numerorum saepenumero opus sit, ne, quae non sunt, pro certis assumamus. Celeb. Lambert in diss. iam supra laudata \emph{Acta Erudit.} 1769 \emph{p.} 127 huius propositionis mentionem facit sed demonstrationis ne necessitatem quidem attigit. Nemo vero demonstrationem tentavit praeter summum Eulerum, \emph{Comment. nov. Ac. Petrop. T.} XVIII ad annum 1773, \emph{Demonstrationes circa residua ex divisione potestatum per numeros primos resultantia} p. 85 seqq. vid. imprimis art. 37 ubi de demonstrationis necessitate fusius locutus est. At demonstratio quam Vir sagacissimus exhibuit duos defectus habet. Alterum quod art. 31 et sqq. tacite supponit, congruentiam $x^n\equiv1$ (translatis ratiociniis illic adhibitis in nostra signa) revera $n$ radices diversas habere, quamquam ante nihil aliud fuerit demonstratum quam quod \emph{plures} habere nequeat; alterum, quod formulam art. 34 per inductionem tantummodo deduxit.

\begin{center}\emph{Radices primitivae, bases, indices.}\end{center}

\Article{57}
Numeros ad exponentem $p-1$ pertinentes \emph{radices primitivas} cum ill. Eulero vocabimus. Si igitur $a$ est radix primitiva, potestatum $a,aa,a^3\ldots a^{p-1}$
% Physical source leaf 052, printed47; same sentence continues.
\SourceLeaf{52}
residua minima omnia erunt diversa; unde facile deducitur, inter haec omnes numeros $1,2,3,\ldots p-1$, qui totidem sunt multitudine quot illa residua minima, reperiri debere, i. e. quemvis numerum per $p$ non divisibilem potestati alicui ipsius $a$ congruum esse. Insignis haec proprietas permagnae est utilitatis, operationesque arithmeticas, ad congruentias pertinentes, haud parum sublevare potest, simili fere modo, ut logarithmorum introductio operationes arithmeticae vulgaris. Radicem aliquam primitivam, $a$, ad lubitum pro \emph{basi} adoptabimus, ad quam omnes numeros per $p$ non divisibiles referemus, et si fuerit $a^e\equiv b\,(\operatorname{mod.}p)$, $e$ ipsius $b$ \emph{indicem} vocabimus. \emph{Ex. gr.} si pro modulo $19$, radix primitiva $2$ pro basi assumatur respondebunt

\begin{center}
\begingroup\small\setlength{\tabcolsep}{1.5pt}
\begin{tabular}{@{}l*{18}{r}@{}}
numeris & 1. & 2. & 3. & 4. & 5. & 6. & 7. & 8. & 9. & 10. & 11. & 12. & 13. & 14. & 15. & 16. & 17. & 18.\\
indices & 0. & 1. & 13. & 2. & 16. & 14. & 6. & 3. & 8. & 17. & 12. & 15. & 5. & 7. & 11. & 4. & 10. & 9.
\end{tabular}
\endgroup
\end{center}

Ceterum patet, manente basi, cuique numero plures indices convenire, sed hos omnes secundum modulum $p-1$ fore congruos; quamobrem quoties de indicibus sermo erit, qui secundum modulum $p-1$ sunt congrui pro aequivalentibus habebuntur, simili modo uti numeri ipsi, quando secundum modulum $p$ sunt congrui, tamquam aequivalentes spectantur.

\begin{center}\emph{Algorithmus indicum.}\end{center}

\Article{58}
Theoremata ad indices pertinentia prorsus analoga sunt iis quae ad logarithmos spectant.

\emph{Index producti e quotcunque factoribus conflati congruus est summae indicum singulorum factorum secundum modulum $p-1$.}

\emph{Index potestatis numeri alicuius congruus est producto ex indice numeri dati in exponentem potestatis, secundum mod. $p-1$.}

Demonstrationes propter facilitatem omittimus.

Hinc perspicitur si tabulam construere velimus ex qua omnium numerorum indices pro modulis diversis desumi possint, ex hac tum omnes numeros modulo maiores, tum omnes compositos omitti posse. Specimen huius modi tabulae ad calcem operis huius adiectum est, \emph{Tab.} I, ubi in prima columna verticali positi sunt numeri primi primorumque potestates a $3$ usque ad $97$, qui tamquam moduli sunt spectandi, iuxta hos singulos numeri pro basi assumti; tum sequuntur indices numerorum primorum successivorum, quorum quini semper per parvulum
% Physical source seam052--053: outgoing in- / incoming tervallum.
intervallum\SourceLeaf{53} sunt disiuncti, eodemque ordine supra dispositi sunt numeri primi; ita ut quis index numero primo dato secundum modulum datum respondeat, facile tutoque inveniri possit.

Ita \emph{ex. gr.} si $p=67$ index numeri $60$, assumto $12$ pro basi erit
\[\equiv2\operatorname{Ind.}2+\operatorname{Ind.}3+\operatorname{Ind.}5\ (\operatorname{mod.}66)\equiv58+9+39\equiv40\]

\Article{59}
\emph{Index valoris cuiuscunque expressionis $\frac{a}{b}(\mathit{mod.}p.),$ (art. 31) congruus est secundum modulum $p-1$ differentiae indicum numeratoris $a$ et denominatoris $b$, siquidem numeri $a,b$ per $p$ non sunt divisibiles.}

Sit enim valor quicunque $c$: eritque $bc\equiv a\,(\operatorname{mod.}p)$; hinc
\[\operatorname{Ind.}b+\operatorname{Ind.}c\equiv\operatorname{Ind.}a\,(\operatorname{mod.}p-1)\]
adeoque
\[\operatorname{Ind.}c\equiv\operatorname{Ind.}a-\operatorname{Ind.}b\]

Si itaque tabula habetur, ex qua index cuique numero respondens pro quovis modulo primo, aliaque ex qua numerus ad indicem datum pertinens derivari possit, omnes congruentiae primi gradus facillimo negotio solvi poterunt, quoniam omnes reduci possunt ad tales, quarum modulus est numerus primus (art. 30).

\emph{E. g.} proposita congruentia
\[29x+7\equiv0\,(\operatorname{mod.}47)\quad\text{erit}\quad x\equiv\frac{-7}{29}(\operatorname{mod.}47)\]
Hinc
\[\operatorname{Ind.}x\equiv\operatorname{Ind.}-7-\operatorname{Ind.}29\equiv\operatorname{Ind.}40-\operatorname{Ind.}29\equiv15-43\equiv18(\operatorname{mod.}46)\]
At numerus cuius index $18$ invenitur $3$. Quare $x\equiv3\,(\operatorname{mod.}47)$. --- Tabulam secundam quidem non adiecimus: at huius vice alia defungi poterit uti Sect. VI ostendemus.

\begin{center}\emph{De radicibus congruentiae $x^n\equiv A$.}\end{center}

\Article{60}
Simili modo ut art. 31 radices congruentiarum primi gradus designavimus. in sequentibus etiam congruentiarum purarum altiorum graduum radices per signum exhibebimus. Uti scilicet ${}^{n}\!\surd A$ nihil aliud significat quam radicem aequationis $x^n=A$, ita apposito modulo per ${}^{n}\!\surd A\,(\operatorname{mod.}p)$ denotabitur radix quaecunque congruentiae $x^n\equiv A\,(\operatorname{mod.}p)$. Hanc expressionem ${}^{n}\!\surd A\,(\operatorname{mod.}p)$ tot valores habere dicemus, quot habet secundum $p$ incongruos, omnes enim qui secundum $p$
% Physical source leaf054, printed49; same sentence continues.
\SourceLeaf{54}
sunt congrui tamquam aequivalentes spectandi (art. 26). Ceterum patet, si $A,B$ secundum $p$ fuerint congrui, expressiones ${}^{n}\!\surd A,{}^{n}\!\surd B\,(\operatorname{mod.}p)$ aequivalentes fore.

Iam si ponitur ${}^{n}\!\surd A\equiv x\,(\operatorname{mod.}p)$, erit $n\operatorname{Ind.}x\equiv\operatorname{Ind.}A\,(\operatorname{mod.}p-1)$. Ex hac congruentia deducuntur ad praecepta sectionis praec. valores ipsius $\operatorname{Ind.}x$ atque ex his valores respondentes ipsius $x$. Facile vero perspicitur $x$ habere totidem valores, quot radices congruentia $n\operatorname{Ind.}x\equiv\operatorname{Ind.}A\,(\operatorname{mod.}p-1)$. Manifesto igitur ${}^{n}\!\surd A$ unum tantummodo valorem habebit quando $n$ ad $p-1$ est primus; quando vero numeri $n,p-1$ divisorem communem habent $\delta$, atque hic est maximus, $\operatorname{Ind.}x$ habebit $\delta$ valores incongruos secundum $p-1$, adeoque ${}^{n}\!\surd A$ totidem valores incongruos secundum $p$, siquidem $\operatorname{Ind.}A$ per $\delta$ est divisibilis. Qua conditione deficiente ${}^{n}\!\surd A$ nullum valorem realem habebit.

\emph{Exemplum.} Quaeruntur valores expressionis ${}^{15}\!\surd11\,(\operatorname{mod.}19)$. Solvi itaque debet congruentia $15\operatorname{Ind.}x\equiv\operatorname{Ind.}11\equiv6\,(\operatorname{mod.}18)$, invenienturque tres valores ipsius $\operatorname{Ind.}x\equiv4,10,16\,(\operatorname{mod.}18)$. His vero respondent valores ipsius $x$, $6,9,4$.

\Article{61}
Quantumvis expedita sit methodus haec quando tabulae necessariae adsunt, debemus tamen non oblivisci, indirectam eam esse. Operae igitur pretium erit inquirere quantum methodi directae polleant: trademusque hic ea quae ex praecedentibus hauriri possunt: alia, quae considerationes reconditiores postulant, ad sectionem VIII reservantes. Initium facimus a casu simplicissimo, ubi $A=1$, sive ubi radices congruentiae $x^n\equiv1\,(\operatorname{mod.}p)$ quaeruntur. Hic itaque, assumta radice quacunque primitiva pro basi, debet esse $n\operatorname{Ind.}x\equiv0\,(\operatorname{mod.}p-1)$. Quae congruentia, quando $n$ ad $p-1$ est primus, unam tantummodo radicem habebit, scilicet $\operatorname{Ind.}x\equiv0\,(\operatorname{mod.}p-1)$: quare \emph{in hocce casu} ${}^{n}\!\surd1\,(\operatorname{mod.}p)$ unicum valorem habet, scilicet $\equiv1$. Quando autem numeri $n,p-1$ habent divisorem communem (maximum) $\delta$, congruentiae $n\operatorname{Ind.}x\equiv0\,(\operatorname{mod.}p-1)$ solutio completa erit $\operatorname{Ind.}x\equiv0\,(\operatorname{mod.}\frac{p-1}{\delta})$ (V. art. 29), \emph{i. e.} $\operatorname{Ind.}x$ secundum modulum $p-1$ alicui ex his numeris
\[0,\frac{p-1}{\delta}.\ \frac{2(p-1)}{\delta},\frac{3(p-1)}{\delta},\GaussFourDots\frac{(\delta-1)(p-1)}{\delta}\]
congruus esse debebit, sive $\delta$ valores secundum modulum $p-1$ incongruos habebit; quare etiam $x$ in hocce casu $\delta$ valores diversos (secundum modulum
